\begin{abstract}
Experience memory---banks of distilled episodes retrieved into an agent's context---is widely reported to improve LLM agents, and its suspected failure mode of premature commitment to remembered answers is a growing concern. We pre-registered a 33-arm, $\sim$29{,}500-episode dissection of an ExpeL-style memory bank on three multi-hop QA benchmarks: the design was frozen by hash before any test-set episode, four numeric go/no-go gate rounds were passed with thresholds never modified, and every deviation is logged. For a strong 27B agent, every registered memory effect vanishes under controls: the primary paired endpoint is formally \emph{equivalent} to no-memory (paired $\Delta$EM $+0.6$pt; TOST 90\% CI $[-1.2,+2.5]$ within $\pm4$pt), six content manipulations (answer redaction, information swap, answer-only, decoys, capacity) are individually null, and the hypothesized search-truncation mechanism is absent---memory \emph{increases} search. What dominates instead is the instruction slot: a one-line evidence-citation instruction yields $+4$--$6$pt (paired exact McNemar $p\le4\times10^{-7}$, Holm-corrected), hard enforcement adds nothing over merely asking, and instructions evolved by GEPA- and MIPROv2-style optimizers (350-rollout budgets, frontier reflection) converge to the hand-written instruction's ceiling---their minibatch gains are scorer adaptation, as a frontier audit of the EM metric confirms (15/15 sampled ``errors'' semantically correct). All of this is agent-relative: for a weaker 120B-MoE agent the structure inverts---memory helps ($+7.7$pt, $p{=}.0018$), value flows through stored answers (redaction $-5.0$pt, $p{=}.040$), and enforcement beats instruction ($+3.3$pt gap). Claims about agent memory, we conclude, are claims about specific agents. The cheapest baseline---one instruction, no memory---was the strongest configuration we found for the strong agent, and pre-registration is what made the null visible.
\end{abstract}
